%% =============================================================================
%%  Etapa 1 — Traduções
%% =============================================================================
\chapter{Traduções}\label{cap:traducoes}

\begin{objetivo}
Ler o trecho em onze traduções portuguesas, notar onde elas concordam e onde
divergem, e transformar cada divergência em pergunta para as etapas seguintes.
Nenhuma tradução é \enquote{a certa}: cada uma resolve de um jeito problemas que
estão no texto original.
\end{objetivo}

\section{Como ler várias traduções}

As traduções diferem por três motivos principais. Primeiro, pela \emph{base
textual}: a Almeida Revista e Corrigida segue o \emph{Textus Receptus}; a
maioria das versões modernas segue o texto crítico (Nestle-Aland/UBS ou
Westcott-Hort); a Bíblia Peshitta traduz o siríaco. Segundo, pela
\emph{filosofia de tradução}: umas buscam correspondência formal (a forma da
frase original), outras equivalência funcional (o efeito no leitor de hoje)
\parencite{feestuart2014}. Terceiro, pela \emph{linguagem}: \enquote{vós}
e \enquote{vocês}, vocabulário erudito ou coloquial.

\begin{nota}[Dica de leitura]
Leia primeiro todas as traduções do versículo 20, depois todas do 21, e assim
por diante. Sublinhe à mão as palavras que mudam de uma para outra. As
caixas de base textual informam o texto hebraico e grego que cada editora
declara ter usado.
\end{nota}

%------------------------------------------------------------------------------
\section{As traduções}

\begin{traducao}{TNM}{Tradução do Novo Mundo da Bíblia Sagrada}{2015}
\vs{20}E ele ergueu os olhos para os seus discípulos e disse:\\
“Felizes são vocês, pobres, pois o Reino de Deus é de vocês.\\
\vs{21}“Felizes são vocês que agora têm fome, pois serão saciados.\\
“Felizes são vocês que agora choram, pois vão rir.\\
\vs{22}“Felizes serão vocês sempre que os homens os odiarem, e sempre que os
excluírem, os insultarem e declararem maldito o seu nome, por causa do Filho do
Homem. \vs{23}Alegrem-se quando isso acontecer e pulem de alegria, porque a sua
recompensa é grande no céu, pois essas são as mesmas coisas que os antepassados
deles faziam aos profetas.\\
\vs{24}“Mas, ai de vocês, ricos, pois já receberam todo o seu conforto.\\
\vs{25}“Ai de vocês que agora estão saciados, pois passarão fome.\\
“Ai de vocês que agora riem, pois lamentarão e chorarão.\\
\vs{26}“Ai de vocês, sempre que todos os homens falarem bem de vocês, pois é isso
que os antepassados deles fizeram aos falsos profetas.
\basetextual{\emph{Biblia Hebraica} (Kittel); consultadas a \emph{Stuttgartensia}
  e a \emph{Quinta}}{Westcott-Hort; consultados Nestle-Aland e UBS}
\end{traducao}

\begin{traducao}{TNM}{Tradução do Novo Mundo da Bíblia Sagrada}{1986}
\vs{20}E ele ergueu os olhos para os seus discípulos e começou a dizer:\\
“Felizes sois vós, pobres, porque vosso é o reino de Deus.\\
\vs{21}“Felizes sois vós os que agora tendes fome, porque sereis saciados.\\
“Felizes sois vós os que agora chorais, porque haveis de rir.\\
\vs{22}“Felizes sois sempre que os homens vos odiarem, e sempre que vos
excluírem, e vos vituperarem, e lançarem fora o vosso nome, como iníquo, por
causa do Filho do homem. \vs{23}Alegrai-vos naquele dia e pulai, pois, eis que a
vossa recompensa é grande no céu, porque estas são as mesmas coisas que os
antepassados deles costumavam fazer aos profetas.\\
\vs{24}“Mas, ai de vós, ricos, porque já tendes plenamente a vossa consolação.\\
\vs{25}“Ai de vós os que agora estais saciados, porque passareis fome.\\
“Ai, vós que agora rides, porque pranteareis e chorareis.\\
\vs{26}“Ai, sempre que todos os homens falarem bem de vós, porque coisas como
essas são as que os antepassados deles fizeram aos falsos profetas.
\basetextual{\emph{Biblia Hebraica} (Kittel); consultadas a \emph{Stuttgartensia}
  e a \emph{Quinta}}{Westcott-Hort; consultados Nestle-Aland e UBS}
\end{traducao}

\begin{traducao}{TB}{Tradução Brasileira}{2010}
\vs{20}Olhando para seus discípulos, começou a dizer:\\
Bem-aventurados vós os pobres, porque vosso é o reino de Deus.\\
\vs{21}Bem-aventurados vós que agora tendes fome, porque sereis fartos.\\
Bem-aventurados vós que agora chorais, porque vos rireis.\\
\vs{22}Bem-aventurados sois quando os homens vos odiarem e quando vos
expulsarem da sua companhia, vos ultrajarem e rejeitarem o vosso nome como
indigno, por causa do Filho do Homem. \vs{23}Regozijai-vos naquele dia e
exultai, porque grande é o vosso galardão no céu; pois assim seus pais
trataram aos profetas.\\
\vs{24}Mas ai de vós que sois ricos! Porque já recebestes a vossa consolação.\\
\vs{25}Ai de vós, os que agora estais fartos! Porque tereis fome.\\
Ai de vós, os que agora rides! Porque haveis de lamentar e chorar.\\
\vs{26}Ai de vós, quando todos vos louvarem! porque assim seus pais trataram
aos falsos profetas.
\basetextual{Bíblia hebraica de Letteris}{Nestle-Aland}
\end{traducao}

\begin{traducao}{RC}{Almeida Revista e Corrigida}{1969}
\vs{20}E, levantando ele os olhos para os seus discípulos, dizia:\\
Bem-aventurados vós, os pobres, porque vosso é o reino de Deus.\\
\vs{21}Bem-aventurados vós, que agora tendes fome, porque sereis fartos.\\
Bem-aventurados vós, que agora chorais, porque haveis de rir.\\
\vs{22}Bem-aventurados sereis quando os homens vos aborrecerem e quando vos
separarem, e vos injuriarem, e rejeitarem o vosso nome como mau, por causa do
Filho do homem. \vs{23}Folgai nesse dia, exultai; porque, eis que é grande o
vosso galardão no céu, pois assim faziam os seus pais aos profetas.\\
\vs{24}Mas ai de vós ricos! porque já tendes a vossa consolação.\\
\vs{25}Ai de vós, os que estais fartos, porque tereis fome.\\
Ai de vós, os que agora rides, porque vos lamentareis e chorareis.\\
\vs{26}Ai de vós quando todos os homens de vós disserem bem, porque assim
faziam seus pais aos falsos profetas.
\basetextual{Texto Massorético}{\emph{Textus Receptus}}
\end{traducao}

\begin{traducao}{ARC}{Almeida Revista e Corrigida}{2009}
\vs{20}E, levantando ele os olhos para os seus discípulos, dizia:\\
Bem-aventurados vós, os pobres, porque vosso é o Reino de Deus.\\
\vs{21}Bem-aventurados vós, que agora tendes fome, porque sereis fartos.\\
Bem-aventurados vós, que agora chorais, porque haveis de rir.\\
\vs{22}Bem-aventurados sereis quando os homens vos aborrecerem, e quando vos
separarem, e vos injuriarem, e rejeitarem o vosso nome como mau, por causa do
Filho do Homem. \vs{23}Folgai nesse dia, exultai, porque é grande o vosso
galardão no céu, pois assim faziam os seus pais aos profetas.\\
\vs{24}Mas ai de vós, ricos! Porque já tendes a vossa consolação.\\
\vs{25}Ai de vós, os que estais fartos, porque tereis fome!\\
Ai de vós, os que agora rides, porque vos lamentareis e chorareis!\\
\vs{26}Ai de vós quando todos os homens falarem bem de vós, porque assim faziam
seus pais aos falsos profetas!
\basetextual{Texto Massorético}{\emph{Textus Receptus}}
\end{traducao}

\begin{traducao}{ARA}{Almeida Revista e Atualizada}{1993}
\vs{20}Então, olhando ele para os seus discípulos, disse-lhes:\\
Bem-aventurados vós, os pobres, porque vosso é o reino de Deus.\\
\vs{21}Bem-aventurados vós, os que agora tendes fome, porque sereis fartos.\\
Bem-aventurados vós, os que agora chorais, porque haveis de rir.\\
\vs{22}Bem-aventurados sois quando os homens vos odiarem e quando vos
expulsarem da sua companhia, vos injuriarem e rejeitarem o vosso nome como
indigno, por causa do Filho do Homem. \vs{23}Regozijai-vos naquele dia e
exultai, porque grande é o vosso galardão no céu; pois dessa forma procederam
seus pais com os profetas.\\
\vs{24}Mas ai de vós, os ricos! Porque tendes a vossa consolação.\\
\vs{25}Ai de vós, os que estais agora fartos! Porque vireis a ter fome.\\
Ai de vós, os que agora rides! Porque haveis de lamentar e chorar.\\
\vs{26}Ai de vós, quando todos vos louvarem! Porque assim procederam seus pais
com os falsos profetas.
\basetextual{\emph{Biblia Hebraica Stuttgartensia}}{Nestle-Aland}
\end{traducao}

\begin{traducao}{NAA}{Nova Almeida Atualizada}{2017}
\vs{20}Então, olhando para os seus discípulos, Jesus lhes disse:\\
— Bem-aventurados são vocês, os pobres,\\
\quad porque o Reino de Deus é de vocês.\\
\vs{21}— Bem-aventurados são vocês que agora têm fome,\\
\quad porque serão saciados.\\
— Bem-aventurados são vocês que agora choram,\\
\quad porque vocês hão de rir.\\
\vs{22}— Bem-aventurados são vocês quando as pessoas os odiarem, expulsarem da
sua companhia, insultarem e rejeitarem o nome de vocês como indigno, por causa
do Filho do Homem. \vs{23}Alegrem-se naquele dia e exultem, porque grande é a
recompensa de vocês no céu; porque os pais dessas pessoas fizeram o mesmo com
os profetas.\\
\vs{24}Mas ai de vocês, os ricos,\\
\quad porque vocês já receberam a consolação!\\
\vs{25}Ai de vocês que agora estão fartos,\\
\quad porque vocês vão passar fome!\\
— Ai de vocês que agora estão rindo,\\
\quad porque vocês vão lamentar e chorar!\\
\vs{26}Ai de vocês, quando todos os elogiarem, porque os pais dessas pessoas
fizeram o mesmo com os falsos profetas!
\basetextual{\emph{Biblia Hebraica Stuttgartensia}}{UBS 5.ª ed. / Nestle-Aland 28.ª ed.}
\end{traducao}

\begin{traducao}{NTLH}{Nova Tradução na Linguagem de Hoje}{2000}
\vs{20}Jesus olhou para os seus discípulos e disse:\\
— Felizes são vocês, os pobres,\\
\quad pois o Reino de Deus é de vocês.\\
\vs{21}— Felizes são vocês que agora têm fome,\\
\quad pois vão ter fartura.\\
— Felizes são vocês que agora choram,\\
\quad pois vão rir.\\
\vs{22}— Felizes são vocês quando os odiarem, rejeitarem, insultarem e disserem
que vocês são maus por serem seguidores do Filho do Homem. \vs{23}Fiquem
felizes e muito alegres quando isso acontecer, pois uma grande recompensa está
guardada no céu para vocês. Pois os antepassados dessas pessoas fizeram essas
mesmas coisas com os profetas.\\
\vs{24}— Mas ai de vocês que agora são ricos,\\
\quad pois já tiveram a sua vida boa.\\
\vs{25}— Ai de vocês que agora têm tudo,\\
\quad pois vão passar fome.\\
— Ai de vocês que agora estão rindo,\\
\quad pois vão chorar e se lamentar.\\
\vs{26}— Ai de vocês quando todos os elogiarem, pois os antepassados dessas
pessoas também elogiaram os falsos profetas.
\basetextual{não informada no material original — \textbf{a conferir} na introdução da edição}%
  {não informada no material original — \textbf{a conferir}}
\end{traducao}

\begin{traducao}{AS21}{Almeida Século 21}{2013}
\vs{20}Então, olhando para os discípulos, disse:\\
Bem-aventurados sois vós, os pobres, porque o reino de Deus é vosso.\\
\vs{21}Bem-aventurados sois vós que agora tendes fome, porque ficareis satisfeitos.\\
Bem-aventurados sois vós que agora chorais, porque havereis de rir.\\
\vs{22}Bem-aventurados sereis quando os homens vos odiarem, quando vos
expulsarem da companhia deles, vos insultarem e rejeitarem o vosso nome como
indigno, por causa do Filho do homem. \vs{23}Alegrai-vos nesse dia e exultai,
porque a vossa recompensa no céu é grande; pois os antepassados deles fizeram o
mesmo com os profetas.\\
\vs{24}Mas ai de vós que sois ricos, porque já recebestes a vossa consolação.\\
\vs{25}Ai de vós que agora estais satisfeitos, porque passareis fome.\\
Ai de vós que agora estais rindo, porque vos lamentareis e chorareis.\\
\vs{26}Ai de vós, quando todos vos elogiarem, porque os antepassados deles
fizeram o mesmo com os falsos profetas.
\basetextual{\emph{Biblia Hebraica Stuttgartensia}}{UBS 4.ª ed. / Nestle-Aland 27.ª ed.}
\end{traducao}

\begin{traducao}{NVT}{Nova Versão Transformadora}{2016}
\vs{20}Então Jesus se voltou para seus discípulos e disse:\\
“Felizes são vocês, pobres,\\
\quad pois o reino de Deus lhes pertence.\\
\vs{21}Felizes são vocês que agora estão famintos,\\
\quad pois serão saciados.\\
Felizes são vocês que agora choram,\\
\quad pois no devido tempo rirão.\\
\vs{22}Felizes são vocês quando os odiarem e os excluírem, quando zombarem de
vocês e os caluniarem como se fossem maus porque seguem o Filho do Homem.
\vs{23}Quando isso acontecer, alegrem-se e exultem, porque uma grande
recompensa os espera no céu. E lembrem-se de que os antepassados deles
trataram os profetas da mesma forma.\\
\vs{24}“Que aflição espera vocês, ricos,\\
\quad pois já receberam sua consolação!\\
\vs{25}Que aflição espera vocês que agora têm fartura,\\
\quad pois um terrível tempo de fome os espera!\\
Que aflição espera vocês que agora riem,\\
\quad pois em breve seu riso se transformará em lamento e tristeza!\\
\vs{26}Que aflição espera vocês que são elogiados por todos,\\
\quad pois os antepassados deles também elogiaram falsos profetas!”
\basetextual{\emph{Biblia Hebraica Stuttgartensia}}{Nestle-Aland}
\end{traducao}

\begin{traducao}{BJ}{Bíblia de Jerusalém}{2002}
\vs{20}Erguendo então os olhos para os seus discípulos, dizia:\\
“Felizes vós, os pobres, porque vosso é o Reino de Deus.\\
\vs{21}Felizes vós, que agora tendes fome, porque sereis saciados.\\
Felizes vós, que agora chorais, porque haveis de rir.\\
\vs{22}Felizes sereis quando os homens vos odiarem, quando vos rejeitarem,
insultarem e proscreverem vosso nome como infame, por causa do Filho do Homem.
\vs{23}Alegrai-vos naquele dia e exultai, porque no céu será grande a vossa
recompensa; pois do mesmo modo seus pais tratavam os profetas.\\
\vs{24}Mas, ai de vós, ricos, porque já tendes a vossa consolação!\\
\vs{25}Ai de vós, que agora estais saciados, porque tereis fome!\\
Ai de vós, que agora rides, porque conhecereis o luto e as lágrimas!\\
\vs{26}Ai de vós, quando todos vos bendisserem, pois do mesmo modo seus pais
tratavam os falsos profetas.
\basetextual{Texto Massorético, seguindo as opções da \emph{Bible de Jérusalem}}%
  {texto crítico grego, seguindo as opções da \emph{Bible de Jérusalem} —
   \textbf{a conferir} (o material original repetia aqui \enquote{Texto Massorético})}
\end{traducao}

\begin{traducao}{BP}{Bíblia Peshitta}{2019}
\vs{20}Alçando então seus olhos aos seus discípulos, disse:\\
Felizes sois vós pobres, porque vosso é o reino de Deus;\\
\vs{21}felizes os que agora têm fome, porque vos saciareis;\\
felizes os que agora choram, porque sorrireis.\\
\vs{22}Felizes sereis quando os homens vos aborreçam, e se apartem de vós e vos
insultarem e rejeitem vosso nome como pessoas más, por causa do Filho do
Homem. \vs{23}Regozijai-vos nesse dia e dançai de júbilo, porque vossa
recompensa é grande no Céu, pois deste modo tratavam os pais deles aos
profetas.\\
\vs{24}Mas ai de vós ricos! Porque já recebestes vossa consolação;\\
\vs{25}ai de vós satisfeitos! Porque padecereis fome.\\
Ai de vós que agora sorríeis! Porque vos lamentareis e chorareis.\\
\vs{26}Ai de vós, quando os homens falarem bem de vós! Pois deste modo tratavam
os pais deles aos falsos profetas.
\basetextual{Peshitta (texto oriental); porções ausentes: \emph{Biblia Hebraica
  Stuttgartensia}}{Peshitta, texto de Urmia}
\end{traducao}

\begin{nota}[Revisão da transcrição]
Em relação ao material original foram corrigidos: \enquote{vooso}
(BP, v.~20) e \enquote{profestas} (BJ, v.~23); a frase repetida
\enquote{quando os homens falarem bem de vós} (BP, v.~26); aspas e asterisco
soltos (AS21, v.~23); \enquote{Sttutgartensia}. Na BP, a quebra
\enquote{ai de vós satisfeitos} foi passada para o v.~25, onde está no grego e
no siríaco. \textbf{A conferir} na edição impressa: a numeração dos
vv.~24–25 e a forma \enquote{sorríeis} (BP); a base textual da BJ e da NTLH.
A menção à \emph{Comma Johanneum} e ao final longo de Marcos, que estava na
caixa da BJ, foi retirada: não se aplica a Lucas e precisa de fonte.
\end{nota}

%------------------------------------------------------------------------------
\section{Comparação dos pontos de divergência}

\begingroup\footnotesize
\begin{ebtabela}{@{}>{\sffamily\bfseries}l >{\raggedright}p{.14\linewidth}
    >{\raggedright}p{.14\linewidth} >{\raggedright}X >{\raggedright}p{.2\linewidth}
    >{\raggedright\arraybackslash}p{.14\linewidth}@{}}
\EBcab{} & \EBcab{\gr{μακάριοι}} & \EBcab{\gr{οὐαί}} & \EBcab{\gr{ἀπέχετε τὴν παράκλησιν} (v.~24)}
  & \EBcab{\gr{ὡς πονηρόν} (v.~22)} & \EBcab{\gr{σκιρτήσατε} (v.~23)}\\\midrule
TNM15 & Felizes & ai & já receberam \emph{todo} o seu conforto & declararem maldito & pulem de alegria\\
TNM86 & Felizes & ai & já tendes \emph{plenamente} a vossa consolação & como iníquo & pulai\\
TB    & Bem-aventurados & ai & já recebestes a vossa consolação & como indigno & exultai\\
RC69  & Bem-aventurados & ai & já tendes a vossa consolação & como mau & exultai\\
ARC   & Bem-aventurados & ai & já tendes a vossa consolação & como mau & exultai\\
ARA   & Bem-aventurados & ai & tendes a vossa consolação & como indigno & exultai\\
NAA   & Bem-aventurados & ai & já receberam a consolação & como indigno & exultem\\
NTLH  & Felizes & ai & já tiveram a sua vida boa & disserem que vocês são maus & fiquem muito alegres\\
AS21  & Bem-aventurados & ai & já recebestes a vossa consolação & como indigno & exultai\\
NVT   & Felizes & Que aflição espera & já receberam sua consolação & caluniarem como se fossem maus & exultem\\
BJ    & Felizes & ai & já tendes a vossa consolação & como infame & exultai\\
BP    & Felizes & ai & já recebestes vossa consolação & como pessoas más & dançai de júbilo\\
\end{ebtabela}
\endgroup

\paragraph{Felizes ou bem-aventurados?} As versões se dividem ao meio.
\enquote{Bem-aventurado} soa religioso e distante; \enquote{feliz} soa como
sentimento. O grego \gr{μακάριος} descreve uma \emph{condição} (de quem está
numa situação favorecida), não um estado de ânimo. Esse ponto volta nas
Etapas~2 e~4.

\paragraph{Receber \enquote{por inteiro}.} Só as duas edições da TNM tornam
explícito que a consolação dos ricos já foi recebida \emph{na totalidade}.
O verbo \gr{ἀπέχω} tem esse matiz técnico, como se verá na Etapa~2.

\paragraph{Onde está o \enquote{agora}?} No v.~25a, RC e ARC dizem só
\enquote{os que estais fartos}; as demais trazem \enquote{agora}. A diferença
vem da base textual e é tratada como variante na Etapa~2.

\paragraph{A Peshitta.} A BP traduz o siríaco e por isso tem soluções
próprias (\enquote{dançai de júbilo}, \enquote{pessoas más}). Na Etapa~5 o
texto siríaco é lido diretamente.

\begin{pergunta}
Que efeito tem, para o leitor, a mudança de \enquote{ai de vós} para
\enquote{que aflição espera vocês} (NVT)? O \enquote{ai} expressa lamento,
ameaça, ou as duas coisas?
\end{pergunta}

\begin{pergunta}
Por que algumas versões dizem \enquote{sereis saciados} e outras
\enquote{sereis fartos}, \enquote{vão ter fartura} ou \enquote{vos
saciareis}? Quem sacia? (Pista: a voz do verbo, na Etapa~2.)
\end{pergunta}

\begin{pergunta}
As felicidades e os ais estão em segunda pessoa (\enquote{vós},
\enquote{vocês}). Isso muda algo em relação a um provérbio em terceira pessoa
(\enquote{felizes os pobres})? Compare com Mateus na Etapa~4.
\end{pergunta}
